\documentclass{article}
\usepackage{enumitem}
\usepackage{amsmath}
\begin{document}

\title{Chapter 14: Respiration in Plants}
\date{}
\maketitle

\begin{enumerate}[label=Q\arabic*.]

\item The net gain of ATP from complete aerobic oxidation of one molecule of glucose (theoretical) is:
\\(A) 2 ATP
\\(B) 8 ATP
\\(C) 36--38 ATP
\\(D) 30 ATP

\item Glycolysis occurs in the:
\\(A) Mitochondrial matrix
\\(B) Cytoplasm (cytosol)
\\(C) Inner mitochondrial membrane
\\(D) Nucleus

\item The end products of anaerobic respiration in yeast are:
\\(A) Lactic acid and CO$_2$
\\(B) Ethanol and CO$_2$
\\(C) Acetyl CoA and H$_2$O
\\(D) Pyruvate and NADH

\item Which molecule acts as the entry point into the Krebs cycle from pyruvate?
\\(A) Oxaloacetate
\\(B) Acetyl CoA
\\(C) Citrate
\\(D) Succinate

\item Assertion: The respiratory quotient (RQ) of carbohydrates is 1.\\
Reason: Equal volumes of CO$_2$ are produced and O$_2$ consumed during complete oxidation of glucose.
\\(A) Both true, Reason is correct explanation
\\(B) Both true, Reason is not correct explanation
\\(C) Assertion true, Reason false
\\(D) Assertion false, Reason true

\item During substrate-level phosphorylation in glycolysis, ATP is produced at which steps?
\\(A) Phosphoglycerate kinase and pyruvate kinase steps
\\(B) Hexokinase and phosphofructokinase steps
\\(C) Glyceraldehyde-3-phosphate dehydrogenase step
\\(D) Enolase and aldolase steps

\item The Krebs cycle is also called the TCA cycle because:
\\(A) It involves transfer of carbon atoms to ATP
\\(B) The first product formed is the tricarboxylic acid citric acid (6C)
\\(C) It occurs in three steps of carboxylation
\\(D) Three molecules of CO$_2$ are produced per turn

\item Oxidative phosphorylation differs from substrate-level phosphorylation in that:
\\(A) Oxidative phosphorylation requires a substrate to directly donate phosphate to ADP
\\(B) Oxidative phosphorylation uses the proton gradient across the inner mitochondrial membrane to drive ATP synthase
\\(C) Oxidative phosphorylation produces less ATP
\\(D) Oxidative phosphorylation occurs in the cytoplasm

\item How many CO$_2$ molecules are released per turn of the Krebs cycle?
\\(A) 1
\\(B) 2
\\(C) 3
\\(D) 6

\item The P/O ratio refers to:
\\(A) Phosphate to oxygen ratio in cell
\\(B) Moles of ATP synthesized per mole of oxygen consumed
\\(C) Proportion of pyruvate oxidized
\\(D) Ratio of phospholipid to oxidase in membrane

\item Which of the following correctly describes the role of NAD$^+$ in cellular respiration?
\\(A) NAD$^+$ directly synthesizes ATP
\\(B) NAD$^+$ accepts electrons and protons to form NADH, which then donates electrons to the ETC
\\(C) NAD$^+$ carries oxygen to tissues
\\(D) NAD$^+$ acts as a structural component of ATP synthase

\item In the electron transport chain, Complex IV (cytochrome c oxidase) directly reduces:
\\(A) NAD$^+$
\\(B) Ubiquinone
\\(C) O$_2$ to H$_2$O
\\(D) Cytochrome b

\item Cyanide poisoning is lethal because it:
\\(A) Inhibits Complex I of the ETC, blocking NADH oxidation
\\(B) Inhibits Complex IV, blocking electron transfer to O$_2$ and stopping ATP synthesis
\\(C) Inhibits ATP synthase directly
\\(D) Inhibits glycolysis by blocking hexokinase

\item Match the following metabolic processes with their locations:\\
1. Glycolysis $\rightarrow$ (a) Cytosol\\
2. Pyruvate oxidation $\rightarrow$ (b) Mitochondrial matrix\\
3. Krebs cycle $\rightarrow$ (c) Mitochondrial matrix\\
4. Oxidative phosphorylation $\rightarrow$ (d) Inner mitochondrial membrane
\\(A) 1-a, 2-b, 3-c, 4-d
\\(B) 1-b, 2-a, 3-d, 4-c
\\(C) 1-c, 2-d, 3-a, 4-b
\\(D) 1-d, 2-c, 3-b, 4-a

\item The RQ for fats is less than 1 (approximately 0.7) because:
\\(A) Fats contain less carbon than carbohydrates
\\(B) Fats are more reduced (higher H:O ratio) and require more O$_2$ for complete oxidation relative to CO$_2$ released
\\(C) Fats produce less energy than carbohydrates
\\(D) Fat oxidation does not produce CO$_2$

\end{enumerate}
\end{document}
